\documentclass[UTF8,zihao=-4,twoside,openany]{ctexbook}

% ==================== 1. 页面几何尺寸（严格对应 Word 页面设置） ====================
\usepackage[
  a4paper,
  top=2.75cm,
  bottom=1.92cm,
  left=3.2cm,
  right=2.2cm,
  headheight=15pt,
  headsep=0.75cm,
  footskip=0.55cm,
  asymmetric % 锁定左侧 3.2cm 装订边距，不随奇偶页翻转
]{geometry}

% 彻底消除 book 类自动插入的多余空白过渡页
\let\cleardoublepage\clearpage

% ==================== 2. 字体与排版引擎底层配置 ====================
\usepackage{fontspec}
\setmainfont{Times New Roman}

% 英文无衬线字体绑定 Arial（本地系统自带）
\IfFontExistsTF{Arial}{%
  \setsansfont{Arial}
  \newfontfamily\arialfont{Arial}
}{%
  \setsansfont{TeX Gyre Heros}
  \newfontfamily\arialfont{TeX Gyre Heros}
}

% 开启伪粗体算法（描边系数 2.5，与 Word 物理墨重完全一致）
\usepackage{xeCJK}
\xeCJKsetup{AutoFakeBold=2.5}

\usepackage{amssymb} % 复选框符号 \square
\usepackage{xurl}    % 允许超长网址任意断开，消除 Overfull \hbox

% 正文基础字号与行距（宋体 12 磅，行距 20 磅）
\renewcommand{\normalsize}{\fontsize{12pt}{20pt}\selectfont}
\normalsize
\setlength{\parindent}{2em}
\setlength{\parskip}{0pt}

% ==================== 3. 标题样式设定（严格依照图示批注） ====================
\ctexset{
  chapter = {
    name = {,},
    number = {\arabic{chapter}},
    format = {\raggedright\bfseries\heiti\fontsize{16pt}{16pt}\selectfont},
    aftername = {\hspace{1em}},
    beforeskip = {24pt},
    afterskip = {18pt},
    fixskip = true
  },
  section = {
    format = {\raggedright\heiti\fontsize{14pt}{14pt}\selectfont},
    aftername = {\hspace{1em}},
    beforeskip = {24pt},
    afterskip = {6pt},
    fixskip = true
  },
  subsection = {
    format = {\raggedright\heiti\fontsize{13pt}{13pt}\selectfont},
    aftername = {\hspace{1em}},
    beforeskip = {12pt},
    afterskip = {6pt},
    fixskip = true
  }
}

% ==================== 4. 页眉页脚配置（彻底剔除模板提示语） ====================
\usepackage{fancyhdr}

% 定义论文真实短标题（用于奇数页页眉联动）
\newcommand{\thesistitle}{基于移动端计算机视觉的抗阻训练速度监控与自主调节研究}

% 前置页样式（摘要、目录：无页眉横线，仅底部居中罗马数字）
\fancypagestyle{frontmatterstyle}{
  \fancyhf{}
  \fancyfoot[C]{\fontsize{10.5pt}{12pt}\selectfont\thepage}
  \renewcommand{\headrulewidth}{0pt}
  \renewcommand{\footrulewidth}{0pt}
}

% 正文页样式（奇数页：真实论文题目；偶数页：学校规范全称）
\fancypagestyle{mainmatterstyle}{
  \fancyhf{}
  \fancyhead[CO]{\fontsize{9pt}{12pt}\selectfont \thesistitle}
  \fancyhead[CE]{\fontsize{9pt}{12pt}\selectfont 广州体育学院 2026届硕士研究生学位论文}
  \fancyfoot[C]{\fontsize{10.5pt}{12pt}\selectfont\thepage}
  \renewcommand{\headrulewidth}{0.75pt}
  \renewcommand{\footrulewidth}{0pt}
}

% ==================== 5. 目录精确样式（黑体 11 磅，单倍行距） ====================
\usepackage{titletoc}
\titlecontents{chapter}[0em]
  {\vspace{6pt}\heiti\fontsize{11pt}{11pt}\selectfont}
  {\thecontentslabel\hspace{1em}}
  {}
  {\titlerule*[0.5pc]{.}\contentspage}

\titlecontents{section}[1.5em]
  {\vspace{6pt}\heiti\fontsize{11pt}{11pt}\selectfont}
  {\thecontentslabel\hspace{1em}}
  {}
  {\titlerule*[0.5pc]{.}\contentspage}

\titlecontents{subsection}[3.0em]
  {\vspace{6pt}\heiti\fontsize{11pt}{11pt}\selectfont}
  {\thecontentslabel\hspace{1em}}
  {}
  {\titlerule*[0.5pc]{.}\contentspage}

% 安全重构目录命令，防范类别码越界崩溃
\makeatletter
\renewcommand{\tableofcontents}{%
  \clearpage
  \addcontentsline{toc}{chapter}{目录}%
  \begin{center}
    {\heiti\bfseries\fontsize{16pt}{16pt}\selectfont 目录}\par
    \vspace{18pt}%
  \end{center}
  \@starttoc{toc}%
}
\makeatother

% ==================== 6. 图表与公式环境配置 ====================
\usepackage{caption}
\DeclareCaptionFont{songti5}{\songti\fontsize{10.5pt}{10.5pt}\selectfont}
\captionsetup[figure]{
  font=songti5,
  name={图},
  labelsep=quad,
  justification=centering,
  skip=6pt
}
\captionsetup[table]{
  font=songti5,
  name={表},
  labelsep=quad,
  justification=centering,
  skip=6pt
}

\usepackage{graphicx}
\usepackage{amsmath}
\usepackage{array}
\usepackage{tabularx}
\usepackage{booktabs}
\usepackage{setspace}
\usepackage{xcolor}
\usepackage[numbers,sort&compress]{natbib}
\usepackage[hidelinks]{hyperref}

% ==================== 7. 自定义下划线与辅助宏 ====================
\newcommand{\fixedunderline}[3][c]{\underline{\makebox[#2][#1]{#3}}}
\newcommand{\dateunderline}{\underline{\makebox[1.2cm]{}}年\underline{\makebox[0.8cm]{}}月\underline{\makebox[0.8cm]{}}日}

% 封面专用 0.8pt 下划线（带 3pt 下沉呼吸感）
\newcommand{\coverline}[3][l]{%
  \makebox[#2][#1]{%
    \rule[-3pt]{0pt}{3pt}%
    \hspace{0.5em}#3%
  }%
  \hspace{-#2}%
  \rule[-3pt]{#2}{0.8pt}%
}

% 校徽智能探测宏
\newcommand{\insertschoollogo}[1]{%
  \IfFileExists{gzsu_logo.png}{\includegraphics[width=#1]{gzsu_logo.png}}{%
  \IfFileExists{gzsu_logo.jpg}{\includegraphics[width=#1]{gzsu_logo.jpg}}{%
  \IfFileExists{logo.png}{\includegraphics[width=#1]{logo.png}}{%
  \IfFileExists{logo.jpg}{\includegraphics[width=#1]{logo.jpg}}{%
  \IfFileExists{logo.pdf}{\includegraphics[width=#1]{logo.pdf}}{%
    \framebox[#1]{\rule{0pt}{1.5cm}\fontsize{10pt}{12pt}\selectfont [请将校徽图片命名为 gzsu\_logo.png 放置于根目录下]}%
  }}}}}%
}

\begin{document}

% ==================== 1. 中文封面 ====================
\begin{titlepage}
  \thispagestyle{empty}
  % 顶部信息栏：通过定宽盒子实现冒号绝对对齐
  \begin{flushleft}
    \fontsize{10.5pt}{18pt}\selectfont
    \makebox[4\ccwd][s]{分类号}：\coverline{3.2cm}{} \hfill 
    \makebox[4\ccwd][s]{学校代码}：\coverline[c]{2.8cm}{10585} \\[0.1cm]
    \makebox[4\ccwd][s]{密级}：\coverline{3.2cm}{} \hfill 
    \makebox[4\ccwd][s]{学号}：\coverline[c]{2.8cm}{******}
  \end{flushleft}

  \vspace{1.0cm}
  \begin{center}
    \insertschoollogo{8.8cm} \\[1.2cm]
    {\heiti\fontsize{26pt}{26pt}\selectfont\ziju{0.35} 硕士学位论文} \\[0.6cm]
    {\heiti\fontsize{14pt}{14pt}\selectfont\ziju{0.15} （专业学位）} \\[1.8cm]
    {\heiti\fontsize{22pt}{28pt}\selectfont 中文题目} \\[2.4cm]

    \fontsize{12pt}{24pt}\selectfont
    \begin{tabular}{r@{\ }l}
      \textbf{培养单位：} & \coverline{6.8cm}{体育教育学院} \\
      \textbf{学位类别：} & \coverline{6.8cm}{体育} \\
      \textbf{专业领域：} & \coverline{6.8cm}{运动训练} \\
      \textbf{研究方向：} & \coverline{6.8cm}{篮球训练} \\
      \textbf{学生姓名：} & \coverline{6.8cm}{XXX} \\
      \textbf{指导教师：} & \coverline{6.8cm}{XXX 教授} \\
      \textbf{第二导师：} & \coverline{6.8cm}{XXX 副教授} \\
    \end{tabular}

    \vfill
    {\songti\fontsize{12pt}{12pt}\selectfont\ziju{0.25} 二〇二*年六月}
  \end{center}
\end{titlepage}

% ==================== 2. 英文封面 ====================
\begin{titlepage}
  \thispagestyle{empty}
  \begin{center}
    \vspace*{0.6cm}
    {\fontsize{14pt}{20pt}\selectfont \textbf{Dissertation Submitted to Guangzhou Sport\\[3pt] University for Master Degree}} \\[1.8cm]

    \IfFileExists{logo_shield.png}{%
      \includegraphics[height=2.3cm]{logo_shield.png}%
    }{%
      \insertschoollogo{2.8cm}%
    } \\[2.0cm]

    {\fontsize{16pt}{22pt}\selectfont \textbf{Thesis/Dissertation Title × × × × ×\\[6pt] × × × × × ×}} \\[3.5cm]
    
    \fontsize{12pt}{22pt}\selectfont
    \begin{tabular}{l@{\hspace{1em}}l}
      Author:     & × × × \\[0.3cm]
      Supervisor: & × × × \\
    \end{tabular}

    \vfill
    {\fontsize{12pt}{12pt}\selectfont \textbf{June, 2026}}
    \vspace*{0.3cm}
  \end{center}
\end{titlepage}

% ==================== 3. 声明与非公开标注页 ====================
\clearpage
\thispagestyle{empty}
\begin{center}
  {\heiti\bfseries\fontsize{16pt}{16pt}\selectfont 广州体育学院学位论文原创性声明}
\end{center}
\vspace{0.4cm}

\noindent 本人郑重声明：本人所呈交的学位论文，是在导师的指导下独立进行研究工作所取得的成果。除文中已经注明引用的内容外，本论文不包含任何其他个人或集体已经发表或撰写过的作品成果，不存在抄袭、篡改、伪造等学术不端行为及买卖论文等造假行为。对本文的研究做出重要贡献的个人和集体，均已在文中以明确方式标明并致谢。本人完全意识到本声明的法律结果由本人承担。

\vspace{0.8cm}
\begin{flushright}
  学生签名：\fixedunderline{3.2cm}{} \qquad \dateunderline \\[0.4cm]
  导师签名：\fixedunderline{3.2cm}{} \qquad \dateunderline
\end{flushright}

\vspace{0.3cm}
\noindent\makebox[\textwidth][c]{\leaders\hbox to 6pt{\hss-\hss}\hfill}
\vspace{0.3cm}

\begin{center}
  {\heiti\bfseries\fontsize{14pt}{14pt}\selectfont 非公开学位论文标注说明}
\end{center}
\vspace{0.2cm}

\noindent 根据广州体育学院有关规定，非公开学位论文须经指导教师同意、作者本人申请和相关部门批准方能标注。未经批准的均为公开学位论文（公开的学位论文本项为空白）。

\vspace{0.3cm}
\noindent
\begin{tabular}{|m{2.3cm}|m{11.8cm}|}
  \hline
  \centering 论文题目 & \\ \hline
  \centering 申请密级 & \quad $\square$限制($\le$2年) \qquad $\square$秘密($\le$10年) \qquad $\square$机密($\le$20年) \\ \hline
  \centering 保密期限 & \qquad\qquad 年\qquad 月\qquad 日 \quad 至 \quad\qquad 年\qquad 月\qquad 日 \\ \hline
  \centering 批准日期 & \qquad\qquad\qquad\qquad 年\qquad 月\qquad 日 \\ \hline
\end{tabular}

\vspace{0.25cm}
\begin{flushright}
  {\fontsize{10.5pt}{12pt}\selectfont 广州体育学院研究生院盖章(有效)\qquad\ }
\end{flushright}
\vspace{-0.1cm}
{\fontsize{9pt}{11pt}\selectfont
\noindent 注：限制$\bigstar$2年(可少于2年); 秘密$\bigstar$10年(可少于10年); 机密$\bigstar$20年(可少于20年)
}

% ==================== 4. 授权书页 ====================
\clearpage
\thispagestyle{empty}
\begin{center}
  {\heiti\bfseries\fontsize{16pt}{16pt}\selectfont 广州体育学院学位论文使用授权书}
\end{center}
\vspace{0.5cm}

\noindent 本人依据广州体育学院关于研究生学位论文收藏和使用的相关规定，同意向学校提交本人的学位论文（纸质版和电子版）。

\vspace{0.3cm}
\noindent 根据《中华人民共和国著作权法》及相关规定，学校有权对本人的学位论文进行使用。具体授权内容包括：

\vspace{0.2cm}
\begin{list}{}{%
  \setlength{\leftmargin}{1.8em}
  \setlength{\itemindent}{-1.2em}
  \setlength{\labelwidth}{1.2em}
  \setlength{\labelsep}{0em}
  \setlength{\parsep}{0pt}
  \setlength{\itemsep}{4pt}
}
  \item[1.\ ] 学校将本人的学位论文编入《广州体育学院硕士学位论文全文数据库》，并提供在图书馆等场所的查阅服务；
  \item[2.\ ] 学校有权将学位论文的全部或部分内容编入相关数据库进行检索，并通过影印、缩印或其他复制方式保存和汇编学位论文，同时根据国家有关规定向指定机构提交；
  \item[3.\ ] 对于非公开的学位论文，待解密后适用本授权书规定。
\end{list}

\vspace{3.5cm}
\begin{flushright}
  学位论文作者签字：\fixedunderline{4.5cm}{} \\[0.8cm]
  \underline{\makebox[1.8cm]{}}年\underline{\makebox[0.9cm]{}}月\underline{\makebox[0.9cm]{}}日\qquad\quad
\end{flushright}

% ==================== 5. 前置内容（摘要、目录） ====================
\frontmatter
\pagestyle{frontmatterstyle}
\pagenumbering{Roman}

% ---------------- 中文摘要 ----------------
\clearpage
\addcontentsline{toc}{chapter}{摘要}
\begin{center}
  {\heiti\bfseries\fontsize{16pt}{16pt}\selectfont 摘要}
  \vspace{18pt}
\end{center}

这部分是论文的摘要

中文摘要是论文内容的简要陈述，一般以第三人称语气撰写，300～1000 字（如遇特殊需要字数可以略多），不加评论和解释，是一篇具有独立性和完整性的短文，能准确反映论文的中心思想，规范的学术用语，逻辑性强、结构严谨，体现出论文的新理论、新方法、新技术等。其内容包括：目的意义、研究方法、研究成果和结论，应与论文等同的主要信息，要突出本论文的创造性成果，不可出现图、表、化学方程式、非公知公用的符号和术语。

关键词是为了便于做文献索引和检索工作而从论文中选取出来用以表示全文主题内容信息的单词或术语，在论文中有明确出处，关键词在摘要内容后另起一行，一般 3～8 个，之间用分号分开。

……

\vspace{0.8cm}
\noindent \textbf{关键词：}术语；最多 7 个汉字；最少 3 组，最多 8 组。

% ---------------- 英文摘要 ----------------
\clearpage
\addcontentsline{toc}{chapter}{ABSTRACT}
\begin{center}
  {\arialfont\bfseries\fontsize{16pt}{16pt}\selectfont ABSTRACT}
  \vspace{18pt}
\end{center}

The following English is not correct.

This part is a summary of my paper.

The Chinese abstract is a brief statement of the content of the paper, which is usually written in third person tone, 300 to 1,000 words (the word count may be slightly extended if special requirements are encountered), without comment and explanation, is a short essay with independence and integrity. The Chinese abstract generally explains the purpose, methods, results and conclusions of the research work, which should include the main information equivalent to the paper, and to highlight the creative results of the paper, and the drawings, tables, chemical equations, nonpublic symbols and terminology should not appear.

Keywords in the abstract content after another line, generally 3$\sim$8, separated by semicolons. The key words are the words or terms that are selected from the paper to express the full text information of the subject content in order to facilitate the document indexing and retrieval work.

\vspace{0.8cm}
\noindent \textbf{Key Words:} Terminology; up to 7 Chinese characters; at least 3 groups, with a maximum of 8 sets.

% ---------------- 目录 ----------------
\tableofcontents

% ==================== 6. 正文主体 ====================
\mainmatter
\pagestyle{mainmatterstyle}
\pagenumbering{arabic}

% 强制正文第一页及各章首页均保持完整双面页眉
\makeatletter
\let\ps@plain\ps@mainmatterstyle
\makeatother

% ---------------- 第 1 章 前言 ----------------
\chapter{前言}

\section{问题的提出}
抗阻训练是维持骨骼肌质量、神经肌肉机能与心代谢健康的核心运动方式，已被美国运动医学学会与世界卫生组织共同列为成年人健康促进的基础训练形式（ACSM, 2011; WHO, 2020）。近年来，随着全民健身战略与《健康中国2030》规划纲要的持续推进，大众群体对抗阻训练的科学认知与参与意愿显著提升，参与规模呈现进一步扩大态势。然而，参与规模的扩大并未伴随专业指导资源的同步供给。当前社会体育指导员队伍的专项技能多集中于传统健身项目（杨楠，2025），而商业化私人教练服务的费用门槛使多数普通训练者难以获得持续的专业指导。这一供需失衡导致大众抗阻训练者普遍处于缺乏专业监督的自我指导状态，需独立完成负荷设定、组数与次数安排以及组间恢复判断等核心训练决策。

自我训练场景中，训练者独立面对两类核心决策挑战：当日训练负荷的科学设定与训练过程中动作技术规范性的持续监控。就负荷设定而言，训练者主要依赖训练前的固定处方或即时主观感受进行判断，缺乏对当日神经肌肉状态的客观量化参照，因而难以准确识别负荷是否偏离目标刺激区间，也难以及时察觉疲劳累积带来的过度训练风险。就动作监控而言，训练者的本体感觉与视角局限使其难以在负荷增大与疲劳累积过程中及时察觉动作代偿与技术失范（Gwynne \& Curran, 2014）。这两类挑战在训练过程中存在交互效应：负荷设定失准会加剧动作技术失范的概率，而动作技术失范则会同时降低训练效益并增大骨骼肌肉系统的损伤风险。要在无专业监督条件下缓解这一双重挑战，训练者需要在训练过程中获取相对可靠的反馈信息以动态调节负荷与动作决策。反馈机制的科学性与可及性因此成为大众抗阻训练科学化的核心议题。

围绕反馈机制的构建，训练科学先后发展出两条方法学路径。第一条路径以主观感知反馈为基础，其代表性工具为主观用力程度量表（Borg, 1982）与储备重复次数评估（repetitions in reserve, RIR）（Zourdos et al., 2016），训练者依据即时用力感与对剩余能力的评估调节训练负荷。这一路径的方法学优势在于低成本与高可及性，其局限在于精度依赖训练经验，缺乏外部客观校准。第二条路径以客观测量反馈为基础，其代表性方法为基于速度的力量训练（velocity-based training, VBT），通过实时监测杠铃向心速度将训练强度锚定在生物力学参数之上（González-Badillo \& Sánchez-Medina, 2010）。这一路径的方法学优势在于测量精度与决策客观性，其局限在于对专业硬件设备的高度依赖。传统 VBT 实现方案依赖线性位置传感器等专业设备，单套价格通常在数千美元级别，这一成本壁垒使其应用场景长期局限于竞技训练与科研实验室。两条路径在较长时期内相对独立地发展，各自积累证据基础，各自服务于不同的应用场景。

近年来，两条路径均出现了实质性变化。在主观反馈侧，基于 RIR 的自主调节在方法学证据上持续积累，包括 RIR 训练与力竭训练在肌肥大适应上的等效性验证（Refalo et al., 2024）。基于这一证据基础，美国运动医学学会于 2026 年发布了距上一版本相隔 17 年的抗阻训练立场声明更新，正式将 RIR 取代传统的 1RM 百分比确立为强度处方的首选指标（Currier et al., 2026）。这一更新意味着主观反馈路径在权威指南层面的方法学地位实质性提升。在客观反馈侧，深度学习姿态估计技术的成熟为 VBT 的低成本实现提供了新的技术路径，基于智能手机摄像头与轻量级姿态估计模型的移动端方案在原理上已具备完成杠铃速度测量任务的技术条件（Bazarevsky et al., 2020）。这一技术演进为客观反馈路径向大众场景的下沉提供了成本与设备可及性上的可能。两条路径的近期变化在时间上高度重合：主观反馈路径在权威指南层面获得方法学地位的提升，客观反馈路径在技术层面获得成本可及性的改善。两条路径的这一同步演进使得一个具有现实意义的科学问题浮现：在规范化 RIR 训练已具备科学性、权威性与可及性的前提下，客观速度反馈是否仍具有值得投入额外成本的适应性增量？

围绕这一核心问题，本研究提出两层递进的追问。第一层追问关乎干预的必要性，即客观速度反馈相较于规范化 RIR 反馈能否在训练执行质量、疲劳控制以及体能适应上产生具有实践意义的增量差异。第二层追问关乎工具的可行性，即若客观反馈确有增量价值，基于计算机视觉的移动端方案能否在真实训练场景中提供接近专业硬件的测量能力与实用价值。为系统回应这两层追问，本研究构建了由三个阶段组成的递进式研究方案：研究一通过测量学效度验证，检验自研移动端 AI 系统与专业硬件在核心运动学参数测量上的一致性，为后续研究提供工具层面的方法学基础；研究二采用两组平行随机对照预试验设计，在真实训练场景下比较基于客观速度反馈与规范化 RIR 反馈的两种自主调节方案在 4 周下肢抗阻训练中的可行性、执行质量与初步适应差异；研究三采用半结构化深度访谈与主题分析，从主观体验层面揭示两种反馈路径下训练者的决策逻辑与工具接受机制。三个阶段依次回应工具效度、干预效果与用户体验三个层面的证据缺口，共同为大众抗阻训练场景下反馈方式的合理选择提供实证参照。

\section{选题的目的意义}

\subsection{选题目的}
本研究围绕 1.1 节提出的两层追问，以青年男性抗阻训练者为研究对象，通过工具效度验证、预试验性随机对照训练干预与半结构化深度访谈相结合的三阶段设计，系统评估基于客观速度反馈与基于规范化 RIR 反馈两条自主调节路径在真实抗阻训练场景中的相对表现。具体研究目的包括以下四个方面：

第一，明确自研移动端 AI 系统在杠铃深蹲平均向心速度与反向运动跳高度测量中的一致性水平、系统偏差方向与绝对误差特征，回答基于计算机视觉的低成本方案在测量学层面能否达到支持真实训练监控应用的效度基础。

第二，评估基于客观速度反馈的抗阻训练方案在真实场景中的招募保留率、训练完成率、速度目标命中率、自动调节规则执行稳定性与不良事件发生情况，对照预设推进标准判断该方案向后续大规模确证性随机对照试验推进的可行性水平。

第三，比较两条反馈路径在训练负荷结构、内部负荷反应与初步体能适应上的差异特征，获取两组在深蹲最大力量、下肢爆发力与训练自我效能上的效应量点估计与置信区间，为后续正式随机对照试验的样本量计算提供先验参数。

第四，从行为决策层面揭示训练者对客观速度反馈与主观 RIR 反馈的接受度差异，识别不同训练经验受试者在自主权让渡、算法信任构建与工具持续使用意愿上的分化特征，为数字化训练工具的用户导向优化提供质性证据。

% 主动分页：彻底根除第 3 页底部的孤挂小标题与“规范/化”、“既/有”断词
\clearpage
\subsection{选题意义}

\noindent \textbf{（一）理论意义}

本研究的理论贡献主要体现在两个维度。第一，本研究将抗阻训练自主调节领域的核心对比框架从“VBT vs. 传统百分比训练”推进至“VBT vs. 规范化 RIR 训练”。既有 VBT 实证研究的对照组普遍设定为传统 \%1RM 处方，其相较于 ACSM 2026 立场声明所推荐的规范化 RIR 训练是否仍能维持适应性优势，目前尚缺乏直接的随机对照证据。本研究通过在同一实验框架下直接比较两条已获科学认可的反馈路径，为抗阻训练自主调节理论在权威指南更新后的方法学再校准提供实证素材。第二，本研究采用了整合测量效度、干预效果与用户体验三个证据层次的评估思路。既有关于数字化训练工具的研究多聚焦于单一维度的测量精度验证，较少将技术工具置于随机对照干预中检验其相较于规范化传统方法的实际增量价值，也较少从用户体验层面追踪工具接受度的分化特征。本研究通过量化过程数据、量化结局数据与质性访谈数据的三角互证，为同类 AI 辅助训练工具的评估方案设计提供了一个具体的方法学案例。

\vspace{0.4cm}
\noindent \textbf{（二）实践意义}

本研究的实践价值按由近及远的顺序体现在三个层次。在直接产出层面，本研究通过对自研移动端 AI 系统在真实训练场景中同步记录数据的分析，揭示了其相较于专业线性位置传感器的一致性水平与系统性偏差特征，为该系统的算法校准与后续版本改进提供了直接的工程参数。在指导实践层面，本研究客观呈现了客观速度反馈与规范化 RIR 反馈在实施成本、执行质量、体能适应与主观体验上的相对表现，为具有类似特征的训练者与体能训练从业者在权衡是否引入客观速度监控时提供了初步的实证参照。在方法学参数层面，本研究作为一项严格遵循预试验性研究规范开展的探索性工作，所报告的招募保留特征、系统可用性瓶颈、各项体能结局的效应量估计与训练执行过程的关键参数，可为后续开展大规模确证性随机对照试验的样本量计算、随访策略设计与干预方案完善提供直接的先验参数。

% ---------------- 第 2 章 文献综述 ----------------
\clearpage
\chapter{文献综述}

\section{反馈在抗阻训练中的作用与类型划分}
抗阻训练的适应效果并非由外部负荷单向决定，而是由训练者能否依据训练过程中获得的反馈信息对负荷进行动态调节所共同塑造。数字化监控辅助训练系统近年来受到广泛关注，其核心价值在于通过通用智能设备为训练者提供更及时、更可执行的反馈信息，从而降低传统专业监控设备的应用门槛。因此，理解此类数字化辅助训练的科学价值，首先需要理解反馈本身在抗阻训练中的作用机制、类型划分及其代表方法的方法学内涵。本节围绕反馈方式这一主线，从训练适应的反馈依赖机制出发，依次论述反馈的生理学与行为学作用机制、反馈的类型划分、两类反馈的代表方法及其相互关系，为后续章节的证据评述奠定统一的概念基础。

\subsection{抗阻训练适应对反馈的依赖}
抗阻训练的适应效果并非由外部负荷单向决定，而是由训练者能否依据训练过程中获得的反馈信息对负荷进行动态调节所共同塑造。数字化监控辅助训练系统近年来受到广泛关注，其核心价值在于通过通用智能设备为训练者提供更及时、更可执行的反馈信息，从而降低传统专业监控设备的应用门槛。因此，理解此类数字化辅助训练的科学价值，首先需要理解反馈本身在抗阻训练中的作用机制、类型划分及其代表方法的方法学内涵。本节围绕反馈方式这一主线，从训练适应的反馈依赖机制出发，依次论述反馈的生理学与行为学作用机制、反馈的类型划分、两类反馈的代表方法及其相互关系，为后续章节的证据评述奠定统一的概念基础。

\subsection{反馈的作用机制与两种类型}
反馈之所以能够改善训练适应，其作用机制可从生理学与行为学两个层面加以理解。在生理学层面，反馈信息的核心功能在于帮助训练者识别当前神经肌肉状态，实现负荷强度与疲劳程度的精准匹配。Sánchez-Medina 与 González-Badillo（2011）的研究证实，杠铃向心速度的下降与神经肌肉疲劳积累和代谢应激水平之间存在稳定的剂量--效应关系，这意味着速度反馈能够为训练者提供关于疲劳累积程度的客观量化参照。同一逻辑也适用于主观反馈，当训练者主观感知到的用力感或剩余重复次数发生变化时，这一变化本身即反映了神经肌肉系统状态的改变。无论反馈的来源是外部测量还是内在感知，其生理学功能都在于将不可见的疲劳过程转化为可操作的调节依据。

在行为学层面，反馈信息通过调节训练者的动作意图与注意焦点影响神经肌肉输出。González-Badillo 等（2014）的研究表明，要求训练者在向心阶段以最大意图速度推动杠铃的神经指令本身即可驱动更高阈值运动单位的募集，而实时速度反馈能够强化这一意图指令的执行质量。Wulf 与 Lewthwaite（2016）提出的 OPTIMAL 运动学习理论从更宏观的视角解释了反馈的行为学效应，该理论指出，增强期望（enhanced expectancies）、自主性支持（autonomy support）与外部注意焦点（external focus of attention）是优化运动表现与学习的三个核心因素。不同类型的反馈在激活这三个因素上呈现不同的侧重：客观速度反馈通过即时的数据呈现同时激活"增强期望"（训练者通过速度数据直观感知自身进步）与"外部注意焦点"（训练者将注意力集中于杠铃速度而非内部肌肉感觉），而主观感知反馈则通过将负荷决策权归还给训练者本人激活"自主性支持"（训练者基于自身评估做出决策）。实时客观反馈通过向训练者呈现直观的达标数据，能够即时强化其对动作执行的控制感，这一行为学机制为客观反馈提升训练者的训练自我效能提供了心理学解释（Wulf \& Lewthwaite, 2016）。这一差异化的行为学作用机制预示了两类反馈可能在训练体验与长期依从性上产生不同影响。

基于反馈信息来源的根本差异，Greig 等（2020）在对自主调节训练术语进行系统梳理时提出，抗阻训练中的反馈可划分为两大类型。第一类是客观测量反馈（objective feedback），其信息载体为外部设备采集的运动学或动力学参数，训练者依据这些参数判断负荷强度与疲劳状态。第二类是主观感知反馈（subjective feedback），其信息载体为训练者自身对用力感与剩余能力的内在评估，无需任何外部设备即可实施。这一划分反映了反馈信息获取的两种根本不同的技术路径：前者依赖外部工具将不可见的力学过程转化为可见的数字信号，后者依赖训练者将内在生理感受转化为可操作的判断。两类反馈在信息来源、精度特征、设备依赖与生态适配性上的系统性差异，直接决定了它们在大众训练场景中的不同适用条件。

\subsection{客观测量反馈的代表方法}
在客观测量反馈的诸多具体形式中，基于速度的力量训练（velocity-based training, VBT）是目前发展最为成熟、证据积累最为充分的方法。VBT 的理论基础源于肌肉收缩中已被充分证实的生物力学规律，即力与速度之间的稳定负相关。González-Badillo 与 Sánchez-Medina（2010）通过对杠铃卧推的系统研究确立，给定相对负荷（\%1RM）与其对应向心速度之间存在高度可预测的线性关系。这一关系的实践意义在于，训练者无需反复进行最大重复重量（1-repetition maximum, 1RM）测试，仅凭杠铃移动速度即可反推当前负荷的相对强度，从而有效规避 1RM 测试本身的日内波动与操作风险。

在此理论基础上，VBT 发展出三类核心操作参数，共同构成从测量速度到调控负荷再到控制疲劳的完整操作链条：
\begin{enumerate}
  \item \textbf{平均向心速度（mean concentric velocity, MCV）}：指杠铃在完整向心阶段的平均移动速度。MCV 是 VBT 最基础的操作指标，其与 \%1RM 之间的线性对应关系使得训练者可以将预设的速度区间直接转化为实时的强度调节指令。例如，若当日目标强度对应的 MCV 区间为 0.55 至 0.65 m/s，训练者可依据实测 MCV 与目标区间的偏差决定是否增加或减轻当前负荷，从而实现相对强度的动态维持。
  \item \textbf{速度损失率（velocity loss, VL\%）}：指同一组内首次与末次重复速度之差占首次速度的百分比。Sánchez-Medina 与 González-Badillo（2011）证实该指标与神经肌肉疲劳积累和代谢应激水平存在稳定的剂量--效应关系，因而可作为组内疲劳的客观替代指标。训练者可依据预设 VL 阈值实现组内疲劳的精准截断，避免主观力竭判断带来的不确定性。Pareja-Blanco 等（2017）进一步揭示，不同 VL 阈值对应方向性不同的适应模式，这一发现将 VL\% 从单纯的疲劳量化指标提升为训练适应的调控工具。
  \item \textbf{最大意图速度（maximum intentional velocity, MIV）}：指训练者在向心阶段以最大努力推动杠铃的神经指令。González-Badillo 等（2014）的研究表明，正是这一意图速度而非实际移动速度驱动了高阈值运动单位的募集。MIV 的方法学意义在于强调 VBT 的有效性不仅依赖于速度的测量精度，也依赖于训练者在每次重复中是否真正执行了最大意图的神经指令。
\end{enumerate}
上述三类参数共同界定了 VBT 作为一套完整方法学的操作内涵，构成了本研究客观反馈干预组的核心方法学基础。

\subsection{主观感知反馈的代表方法}
与 VBT 依赖外部设备采集客观数据形成对照，主观感知反馈路径通过训练者自身的内在评估实现负荷调节，其代表方法的方法学发展经历了从 RPE 到 RIR 的关键演进。Borg（1982）最初提出的主观用力程度量表（rating of perceived exertion, RPE）通过整合呼吸困难、肌肉疲劳与主观努力程度等多维信号对运动强度进行整体量化。RPE 作为主观反馈的开创性工具，最初主要应用于有氧运动情境，后逐步拓展至抗阻训练。

然而，传统 RPE 在抗阻训练场景中存在一个操作性局限，即它给出的是一个笼统的用力程度评分，训练者难以据此直接判断本组还能再做几次或下一组应加多少重量。这一局限的根源在于 RPE 的原始设计假定运动强度是一个连续渐变的过程，而抗阻训练的强度调控更依赖于离散的重复次数与负荷组合决策。为解决这一问题，Zourdos 等（2016）针对抗阻训练的动作特性开发了基于储备重复次数（repetitions in reserve, RIR）的 RPE 量表，将抽象的用力感转化为一个具有直接操作价值的判断，即在维持标准动作技术的前提下本组结束时尚能完成的极限重复次数。

RIR 相较于传统 RPE 的方法学优势体现在三个层面：其一是\textbf{操作直接性}，RIR 不仅是一个感知指标，更是一个负荷调节指令。RIR 为 0 意味着已达技术力竭，RIR 为 2 至 3 意味着保留适度余量，训练者可据此直接决定是否继续加重、维持当前负荷或结束当前组次；其二是\textbf{与动作力学的直接对应性}，RIR 以"还能完成几次"为评估单位，与抗阻训练中"组数、次数、负荷"的三维处方框架形成天然对应，避免了 RPE 数值向具体训练决策转化时的信息损失；其三是\textbf{与客观参数的可桥接性}，Jukic 等（2024）通过对 RIR 与杠铃速度关系的个体化建模证实，个体化的 RIR--速度关系可实现可接受的 RIR 预测准确性，这为主观反馈与客观反馈之间的方法学融合提供了实证基础。上述三方面的方法学优势共同确立了 RIR 作为主观反馈路径当前最具代表性的操作工具地位，也构成了本研究规范化主观反馈干预组的核心方法学基础。

\subsection{两类反馈路径的关系与本章分析框架}
综合上述分析，客观测量反馈与主观感知反馈共享自主调节训练这一上位概念，在逻辑上处于并列而非替代关系。两条路径各有其方法学优势与适用边界，客观反馈以测量精度与决策客观性见长，主观反馈以低成本与高可及性见长。两条路径在较长时期内相对独立地发展，各自积累证据基础，各自服务于不同的应用场景。

两条路径的代表方法（VBT 与 RIR）在方法学发展逻辑上呈现出趋同性。VBT 通过 MCV、VL\%、MIV 三类参数将客观速度信号转化为可操作的负荷调节指令，RIR 通过"还能完成几次"的判断将主观用力感转化为可操作的负荷调节指令，两者最终都指向"如何将某一类反馈信号转化为具体训练决策"这一共同的方法学目标。这一趋同性为两条路径的直接比较提供了逻辑基础，也预示了两条路径在训练适应上可能产生相似的效果。近年来 Naclerio 等的直接对比研究（2025, 2026）初步观察到的"可比"结果，正是这一方法学趋同性在实证层面的反映。

基于上述概念框架，本章后续三节将围绕两条反馈路径的比较展开系统评述。2.2 节梳理客观反馈路径从力--速度关系的实验室奠基到移动端计算机视觉技术推动的应用下沉的完整历史演进，并系统评述 VBT 相较于传统百分比训练的实证证据与技术实现现状。2.3 节梳理主观反馈路径从 RPE 奠基到 RIR 转化再到 ACSM 2026 权威转向的方法学演进，并系统评述 RIR 的效度证据、力竭训练必要性的争议解决以及主观反馈的适用边界与局限。2.4 节作为全章核心，直接比较两条路径的训练效果差异，分析"可比"结果的机制解释，权衡应用层面的边际效益，并在此基础上识别既有文献的关键局限与本研究的定位。

% ---------------- 2.2 客观反馈路径 ----------------
\clearpage
\section{客观反馈路径的方法学演进与实证进展}

\subsection{客观反馈路径的方法学演进}
客观反馈路径的发展跨越近九十年，经历四个关键技术驱动节点：

\noindent \textbf{1. 力--速度关系的理论奠基} \\
Hill（1938）通过离体肌肉收缩实验确立了张力与缩短速度之间的双曲线函数关系（Hill 方程），构成了客观反馈路径最底层的生物力学基石。然而，受制于当时缺乏可在训练现场实时采样的传感装置，该理论在此后半个世纪内主要局限于单关节离体实验室研究，未能转化为可操作的训练负荷调控手段。

\vspace{0.2cm}
\noindent \textbf{2. 负荷--速度关系在训练场景中的验证} \\
进入 21 世纪，商用线性位置传感器（LPT）的出现打破了硬件壁垒。González-Badillo 与 Sánchez-Medina（2010）系统证实了多关节复合动作中相对负荷（\%1RM）与平均向心速度之间的线性负相关，首次实现了在现场训练中"通过速度反推相对强度"，使无损伤风险的动态强度标定成为现实。后续团队进一步验证了该规律在深蹲、卧推等核心复合动作中的普遍稳定性（Sánchez-Medina et al., 2017）。

\vspace{0.2cm}
\noindent \textbf{3. 速度损失指标的方法学确立} \\
Sánchez-Medina 与 González-Badillo（2011）证实了组内速度衰减率（VL\%）与代谢产物积累及神经疲劳之间的剂量--效应关系，确立了客观截断疲劳的量化标准。在此基础上，Pareja-Blanco 等（2017）发现不同速度损失阈值导向不同维度的神经肌肉适应，标志着 VBT 由单一的负荷标定工具升级为训练适应方向的精细化调控体系。

\vspace{0.2cm}
\noindent \textbf{4. 计算机视觉技术的应用引入} \\
针对传统专业设备的高昂成本壁垒，Bazarevsky 等（2020）研发的 BlazePose 模型验证了轻量级姿态估计网络在移动消费级硬件上进行实时运动学追踪的可行性。此后，依托智能手机摄像头的无标记点视觉测速算法迅速迭代，推动了客观速度反馈由专业实验室向大众自我训练场景下沉的技术转向。

\subsection{VBT 相较于传统百分比训练的适应性证据}
\noindent \textbf{1. 伞形综述证据} \\
2025 年发表的一项纳入 AMSTAR-2 方法学评估的伞形综述确认，VBT 相较于传统百分比训练（PBT）在最大力量、爆发力及冲刺速度维度具有一致的高等级证据支持，尤其在设置 10\% 至 20\% 速度损失阈值区间时表现出最优的适应收益稳定性。

\vspace{0.2cm}
\noindent \textbf{2. Meta 分析证据} \\
近五年三项独立的 Meta 分析均支持上述结论。W{\l}odarczyk 等（2021）针对 15 项 RCT 研究的分析证实，10\% 至 20\% VL 阈值方案在爆发力及变向能力上的提升显著优于 PBT；Held 等（2022）的网络 Meta 分析揭示了速度损失阈值与力量素质适应之间的剂量--效应关系；Luo 等（2026）纳入 17 项 RCT 的最新 Meta 分析进一步证实，VBT 在跳跃高度（$\text{SMD}=0.27$，$p<0.05$）与变向速度（$\text{SMD}=0.45$，$p<0.01$）上优于 PBT。

\vspace{0.2cm}
\noindent \textbf{3. 随机对照试验与即时强化效应} \\
Dorrell 等（2020）为期 6 周的干预试验证明，VBT 组在达成同等力量提升的同时，总训练量显著减少，呈现出更优的适应效率；Banyard 等（2020）亦证实了其在下肢爆发力维度的增益。Weakley 等（2020）通过控制试验证明，实时视觉或听觉速度反馈能够通过激励神经募集意图，即时使单次动作的速度与功率输出提升达 10\%。

然而，既有实证研究的对照组几乎全部设定为传统静态 PBT 处方。由于 PBT 本身存在忽视日内能力波动、固定次数无法匹配个体代谢耐受等缺陷，VBT 展现的优势究竟来自其方法学本身，抑或仅仅是对照方案缺陷的反衬，亟待以现代化自主调节方案（RIR）作为新型参照系予以重验。

\subsection{客观反馈的技术实现与效度证据}
\noindent \textbf{1. 线性位置传感器（LPT）} \\
LPT（如 GymAware、T-Force）通过机械缆绳直接位移求导测量速度，不涉及积分误差。Weakley 等（2021）的系统综述表明，GymAware 在平均速度测量上的效度极高（$r=0.90\sim1.00$，$\text{SEE}=0.01\sim0.08\text{ m/s}$，$\text{LOA}=-0.03\sim0.03\text{ m/s}$），被学术界公认为 VBT 测量的效度金标准（Pérez-Castilla et al., 2019）。但单台超过 2000 美元的硬件成本与固定线缆约束，使其难以在大众场景普及。

\vspace{0.2cm}
\noindent \textbf{2. 可穿戴加速度计（IMU）} \\
IMU 方案通过二次积分推算位移与速度。Pérez-Castilla 等（2022）的严谨评估证实，PUSH Band 2.0 虽然信度良好（$\text{CV}=5.81\%\sim13.14\%$），但在中高负荷下的并发效度偏低（$r=0.58\sim0.82$），被明确判定为"高信度但低效度（reliable but invalid）"的典型，因积分漂移与非线性运动误差而难以支撑高精度调控需求。

\vspace{0.2cm}
\noindent \textbf{3. 计算机视觉方案（CV）} \\
无标记点 CV 方案（如 Metric VBT、Vitruve App）借助移动端摄像头完成自动化轨迹捕捉。Balsalobre-Fernández 等（2017）证实了其在理想环境下的准金标准表现。然而，Renner 等（2024）针对真实场景的效度评估暴露出严重缺陷：在 589 次测试中，Metric VBT 漏检率达 8.8\%，MyLift 漏检率达 29.7\%；且杠铃偏转、广角畸变与单侧非对称运动均会导致显著的算法系统误差。

综上，技术方案在"精度、成本、便携度"上呈阶梯互补特征（LPT 精度最高但昂贵；CV 便携且免费，但受环境敏感度与漏检制约）。

% ---------------- 2.3 主观反馈路径 ----------------
\clearpage
\section{主观反馈路径的方法学演进与实证进展}

\subsection{从 RPE 到 RIR 的方法学演进}
Borg（1982）提出的 RPE 量表奠定了感知强度量化的基石，但其连续性用力感评估难以直接映射为抗阻训练离散的组数与次数决策。Zourdos 等（2016）针对抗阻训练特点提出的储备重复次数（RIR）量表，通过精准锚定"距离技术失范还剩几次极限完成能力"，将主观感知直接转化为加减负荷与组内停组的离散操作指令。

\subsection{RIR 训练的效度证据与力竭训练争议}
Jukic 等（2024）证实，个体化的 RIR--速度关系能够达成良好的预测效度。关于抗阻训练是否必须力竭，Grgic 等（2022）的元分析及 Refalo 等（2024）的对照试验证明：保留 1 至 3 次 RIR 的非力竭训练在肌肥大与力量增长上与力竭训练具有等效性，且能够显著减少神经衰竭应激、降低动作形变风险并提升训练依从性。

\subsection{ACSM 2026 立场声明的方法学转向}
美国运动医学学会于 2026 年 4 月正式发布的更新版抗阻训练立场声明（Currier et al., 2026），构成了主观反馈路径发展史上的纲领性里程碑：
\begin{enumerate}
  \item \textbf{证据整合层级}：该声明依托对 137 项系统综述与逾 3 万名健康成年受试者大数据的伞形综合，确立了现代抗阻训练指南。
  \item \textbf{首选强度处方转向}：正式废止了将静态 \%1RM 作为首选处方的传统准则，确立以基于 RIR 的自主调节作为首选强度设定方案（如"保留 2 至 3 次 RIR 完成 3 至 4 组"），将处方基准由静态测试彻底切换为训练动态响应。
  \item \textbf{力竭观念更替}：官方确立了"训练应当接近但不必达到力竭"的指导方针。
\end{enumerate}

\subsection{主观反馈的适用边界与局限}
主观反馈虽然成本为零，但其应用受制于三大边界约束：
\begin{enumerate}
  \item \textbf{经验依赖性}：Hackett 等（2017）证明，初学者的 RIR 预测误差高达 2 至 3 次，而高水平者可稳定控制在 1 次以内；Zourdos 等（2016）证实轻负荷、多次数下主观判断的离散度显著增大。
  \item \textbf{力竭距离敏感性}：Bastos 等（2024）指出，距离力竭越远，受试者对剩余次数的评估偏差越大。Laces 等（2025）进一步发现特殊人群（如老年群体）存在接近力竭时过度高估剩余能力的逆向偏差。
  \item \textbf{定义模糊性}：学术界关于技术力竭、向心力竭与意志力竭的概念混淆，制约了主观反馈标准的绝对同质化。
\end{enumerate}

% ---------------- 2.4 两条反馈路径的直接对比 ----------------
\clearpage
\section{两类反馈路径的直接对比与实际应用权衡}

\subsection{直接对比研究的证据现状}
直接对比 VBT 与 RIR 的实证研究近年方才破题。Naclerio 等在卧推（2025）与深蹲（2026）动作中的急性对照研究表明，客观速度阈值方案与主观 RPE 方案在总负荷量、平均向心速度及疲劳控制等维度均呈现高度"可比性（comparable）"。Orange 等（2022）的系统评价及 2025 年网络元分析进一步证实，当对照组由传统 PBT 替换为现代主观自主调节方案后，VBT 的相对效应优势收窄至小效应（$\text{SMD}=0.28$）。Larsen 等（2021）亦指出两者在促进最大力量方面的长期效果未见统计学差异。

\subsection{"可比结果"的机制解释}
\noindent \textbf{1. 共同的生理驱动核：力竭接近度} \\
无论是客观速度损失截断，抑或主观 RIR 余量控制，两者在神经生物学层面最终执行的是同一调节任务——将运动单位募集维持在目标疲劳区间（Grgic et al., 2022; Jukic et al., 2024），因而生理驱动刺激趋同。

\vspace{0.2cm}
\noindent \textbf{2. 实验周期与效应测度局限} \\
现有直接对比研究多为急性测试或 2 周以内的短周期评估。依据 Rhea（2004）及 Swinton 等（2023）的研究，力量适应效应具有时间依赖性累积特征，短周期内微弱的增量优势易被测量噪声与个体间变异所遮蔽。

\vspace{0.2cm}
\noindent \textbf{3. 行为学激励维度的作用分化} \\
依 OPTIMAL 运动学习理论（Wulf \& Lewthwaite, 2016），VBT 强化外部注意焦点与胜任期望，而 RIR 强化自主决策认同。行为学路径的分歧可能更深远地反映在心理自我效能与依从习惯上，而非单纯的最大力量数值。

\subsection{应用层面的边际效益权衡}
\noindent \textbf{1. 精度增量与经验递减} \\
对经验丰富的高阶训练者，RIR 误差已压缩至 1 次以内，引入外部客观仪器的边际效益微弱；对本体感觉模糊的进阶期大众，外部数据则具备校准其主观偏倚的刚性价值。

\vspace{0.2cm}
\noindent \textbf{2. 成本设备门槛与操作摩擦力} \\
RIR 具备"零成本、零设备"的绝对生态优势；移动端 CV 方案虽省去了专用传感器，但环境布光、支架机位、算法漏检等操作摩擦依然影响大众自我训练的长期依从。

\vspace{0.2cm}
\noindent \textbf{3. 自主权让渡与反馈依赖风险} \\
自决理论（Ryan \& Deci, 2017）表明，自主决策感是维持内在动机的核心。Halperin 等（2018）警示，过度将决策权移交给算法指令可能削弱训练自主性；而引导假说（Winstein \& Schmidt, 1990）进一步指出，高频外部数据依赖可能导致本体感受钝化。

\subsection{既有研究的主要局限}
文献述评揭示出当前领域存在四大关键盲区：
\begin{enumerate}
  \item \textbf{中长期随机对照干预证据缺失}：现有直接对比几乎全为急性单次试验，缺乏 4 周以上中长期适应追踪。
  \item \textbf{"可比性"背后的过程生理机制未明}：缺乏对内部负荷、代偿演进及动作技术损耗的深层追踪。
  \item \textbf{移动端 CV 系统的真实场景生态效度不明}：实验室无干扰条件下的算法准确度无法等同于真实健身房干扰场景下的有效可用度。
  \item \textbf{训练者质性体验维度的研究空白}：严重忽视训练者在真实决策中对算法信任、自主感及心理认知负荷的质性体验。
\end{enumerate}

\subsection{本研究的回应定位}
本研究通过三阶段层层递进的设计回应上述缺口：
\begin{enumerate}
  \item \textbf{研究一（测量学效度验证）}：在标准化与扰动环境下，确立自研移动端 AI 视觉系统与专业 LPT 设备的一致性界限及误差分布。
  \item \textbf{研究二（预试验性随机对照干预）}：采用平行随机对照设计，在 4 周下肢抗阻干预中直接比较移动端 VBT 组与规范化 RIR 组，系统捕获招募保留率、规则执行度、最大力量增量及训练自我效能，为确证性 RCT 提供先验参数。
  \item \textbf{研究三（质性深度访谈）}：运用半结构化访谈与主题分析，穿透量化数字，解析两组训练者在自主权让渡、人机算法信任与长期行为维持上的分化机制。
\end{enumerate}

\subsection{本章小结}
综合梳理可见，主观 RIR 与客观 VBT 均属科学自主调节体系。当评价基准推进至权威的 RIR 方案后，客观速度反馈的价值已非单纯的"力量全面胜出"，而转向"负荷决策质量、疲劳截断边界与体验赋能"等精细增量。图 \ref{fig:ch2-framework} 概括了本研究的理论推演与分析路径。

\begin{figure}[htbp]
  \centering
  \small
  \begin{tabular}{|c|}
    \hline
    \textbf{大众自我训练的核心挑战} \\
    当日负荷设定不确定性 \quad $\longleftrightarrow$ \quad 动作技术与疲劳监控受限 \\
    $\Downarrow$ \\
    \textbf{自主调节训练框架（Autoregulation）} \\
    \begin{tabular}{c@{\hspace{2.5cm}}c}
      \textbf{主观感知反馈路径} & \textbf{客观测量反馈路径} \\
      $\text{RPE} \longrightarrow \text{RIR 量表} \longrightarrow \text{ACSM 2026 首推}$ & $\text{力--速度关系} \longrightarrow \text{LPT} \longrightarrow \text{移动端 CV}$ \\
      优势：零成本、高可及性 & 优势：高精度、客观量化 \\
      局限：强经验依赖、远端误差大 & 局限：高成本设备壁垒、移动端漏检 \\
    \end{tabular} \\
    $\Downarrow$ \\
    \textbf{既有文献核心缺口} \\
    $\text{VBT}$ 相对规范化 $\text{RIR}$ 的直接实证中长期证据稀缺；低成本 $\text{CV}$ 真实生态效度待验 \\
    $\Downarrow$ \\
    \textbf{本研究三阶段回应方案} \\
    研究一：移动端视觉效度验证 \ $\longrightarrow$ \ 研究二：4 周预试验 RCT 对照 \ $\longrightarrow$ \ 研究三：质性决策机制访谈 \\
    \hline
  \end{tabular}
  \vspace{6pt}
  \caption{本研究的概念框架与分析路径}
  \label{fig:ch2-framework}
\end{figure}
% ---------------- 第 3 章 研究设计 ----------------
\clearpage
\chapter{研究设计}
\section{研究对象}
$\times\times\times\times\times\times\times\times\times\times\times\times\times\times\times\times\times\times\times\times\times\times\times\times\times\times\times\times\times\times\times\times\times\times\times\times\times\times\times\times\times$

\section{研究方法}
本部分应当对使用的研究方法做简要的介绍，根据实际情况填写。

$\times\times\times\times\times\times\times\times\times\times\times\times\times\times\times\times\times\times\times\times\times\times\times\times\times\times\times\times\times\times\times\times\times\times\times\times\times\times\times\times\times$

\subsection{调查法}
$\times\times\times\times\times\times\times\times\times\times\times\times\times\times\times\times\times\times\times\times\times\times\times\times\times\times\times\times\times\times\times\times\times\times\times\times\times\times\times\times\times$

\subsection{实验法}
$\times\times\times\times\times\times\times\times\times\times\times\times\times\times\times\times\times\times\times\times\times\times\times\times\times\times\times\times\times\times\times\times\times\times\times\times\times\times\times\times\times$

\subsection{统计分析法}
量化研究，需要详细撰写统计分析方法。$\times\times\times\times\times\times\times\times\times\times\times\times\times\times\times\times\times\times\times\times\times\times\times\times\times\times\times\times\times\times\times\times\times\times\times\times\times\times\times\times\times$

\section{技术路线}
本部分是对论文研究的技术路线与整体结构做一个阐述，就是第一部分主要研究什么，第二部分主要研究什么等等，可以加上研究的技术路线图。

$\times\times\times\times\times\times\times\times\times\times\times\times\times\times\times\times\times\times\times\times\times\times\times\times\times\times\times\times\times\times\times\times\times\times$。

% ---------------- 第 4 章 研究结果与分析 ----------------
\clearpage
\chapter{研究结果与分析}
\section{研究结果1}
当广义控制截面具有式（3）的广义本构关系时，可定义如下的截面示性数
\begin{equation}
  \Phi(\theta) = 
  \begin{cases}
    0, & \text{公式居中，序号顶右} \\
    1, &
  \end{cases}
\end{equation}
$\times\times\times\times\times\times\times\times\times\times\times\times\times\times\times\times\times\times\times\times\times\times\times\times\times\times\times\times\times\times\times\times\times\times\times\times\times\times\times\times\times$，如图 1。

\begin{figure}[htbp]
  \centering
  \insertschoollogo{5.5cm}
  \vspace{6pt}
  \caption{$\times\times\times\times\times\times\times\times$示意图}
  \vspace{12pt}
\end{figure}

$\times\times\times\times\times\times\times\times\times\times\times\times$如表 1。

\begin{table}[htbp]
  \centering
  \caption{杆件状态与非线性构形状态对应关系}
  \vspace{6pt}
  \fontsize{10pt}{14pt}\selectfont
  \begin{tabular*}{\textwidth}{@{\extracolsep{\fill}}ccccccccc@{}}
    \toprule
    $\times\times\times\times\times$ & 1 & 2 & 3 & 4 & 5 & 6 & 7 & 8 \\
    \midrule
    $\times\quad\times\times$ & -- & $\circledast_2$ & $\circledast_3$ & $\circledast_1$ & $\circledast_2,\circledast_3$ & $\circledast_1,\circledast_2$ & $\circledast_1,\circledast_3$ & $\circledast_1,\circledast_2,\circledast_3$ \\
    $\times\quad\times\times\times$ & $\circledast_1,\circledast_2,\circledast_3$ & $\circledast_1,\circledast_3$ & $\circledast_1,\circledast_2$ & $\circledast_2,\circledast_3$ & $\circledast_1$ & $\circledast_3$ & $\circledast_2$ & -- \\
    \bottomrule
  \end{tabular*}
\end{table}

\section{研究结果2}

% ---------------- 第 5 章 讨论 ----------------
\clearpage
\chapter{讨论}
讨论内容一般采用漏斗式结构（从具体到一般，再从一般回到具体）

\section{重申核心发现}
不要直接复制摘要或结论的句子，而是用新的、更概括的语言重申你的主要发现，作为讨论的开始。

\section{解释与比较：与已有文献对话}
这是讨论部分的核心。你需要逐条解释你的主要发现，并将其与前人的研究进行比较。

如果结果与前人一致：说明你的研究验证/支持了前人的结论。

如果结果与前人不同或更新：不要回避差异。大胆提出你的解释（可能是方法不同、样本不同、或者发现了新的特殊情况）。

\section{机理/原因分析：深挖“为什么”}
不要停留在表面描述，要深入探讨这些结果背后的机制、原因或逻辑。

\section{引申意义：理论和实践价值}
阐述研究有什么用处。

理论意义：对现有理论有什么补充、修正或挑战？

实践意义：对现实生活、政策制定、体质健康等有什么指导作用？

\section{研究局限性：坦诚研究的不足}
任何研究都有局限。主动指出局限性不仅不会减分，反而显得你严谨、客观。通常从样本、方法、范围三个角度考虑。

\section{未来展望：为后续研究铺路}
基于上述局限性，或者基于你发现的新问题，提出下一步的研究方向。

% ---------------- 第 6 章 结论与建议 ----------------
\clearpage
\chapter{结论与建议}
\section{结论}
这是对研究结果的最终总结。不要简单重复摘要，也不要复制结果部分的每一条数据，而是要对核心发现进行提炼和升华。

建议使用分点（1., 2., 3.）或分段的形式，清晰罗列。
\begin{enumerate}
  \item $\times\times\times\times\times$
  \item $\times\times\times\times\times$
  \item $\times\times\times\times\times$
\end{enumerate}

\section{建议}
这是“学以致用”的部分。研究结果对于现实世界有什么指导意义。

根据论文实际情况撰写。一般情况下，论文的建议 2--3 条即可，不宜过多，根据本人论文的实际情况填写。

% ==================== 7. 参考文献 ====================
\clearpage
\chapter{参考文献}
{\fontsize{10.5pt}{16pt}\selectfont
{\color{red} 参考以下各类文献的撰写方法，使用顺序编码制编写，从 1 开始排序，是大排序而不是分类排序，以下仅演示各类参考文献的格式及字体字号等。}

\vspace{0.4cm}
\noindent {\color{red} \textbf{1．普通图书}}\\
{[1]} 美国妇产科医师学会．新生儿脑病和脑性瘫痪发病机制与病理生理[M]．段涛，杨慧霞，译．北京：人民卫生出版社，2010:38-39．\\
{[2]} 汪昂．增订本草备要:四卷[M]．刻本．京都：老二酉堂，1881(清光绪七年)．\\
{[3]} 蒋有绪，郭泉水，马娟，等．中国森林群落分类及其群落学特征[M]．北京：科学出版社，1998．\\
{[4]} 侯文顺．高分子物理：高分子材料分析、选择与改性[M/OL]．北京：化学工业出版社，2010：119[2012-11-27]．\url{http://apabi.lib.pku.edu.cn/usp/pku/pub.mvc?pid=book.detail&metaid=m.20111114-HGS-889-0228}．\\
{[5]} CRA WFP RD W，GORMAN M．Future libraries：dreams，madness，\& reality[M]．Chicago：American Library Association，1995．\\
{[6]} International Federation of Library Association and Institutions．Names of persons: national usages for entry in catalogues[M]．3rd ed．London:IFLA International Office for UBC，1977．\\
{[7]} O'BRIEN J A．Introduction to information systems[M]．7th ed．Burr Ridge，III：Irwin，1994．

\vspace{0.4cm}
\noindent {\color{red} \textbf{2．论文集、会议录}}\\
{[1]} 中国职工教育研究会．职工教育研究论文集[G]．北京：人民教育出版社，1985．\\
{[2]} 雷光春．综合湿地管理：综合湿地管理国际研讨会论文集[C]．北京：海洋出版社，2012．\\
{[3]} BABU B V, NAGAR A K, DEEP K, et al. Proceedings of the Second International Conference on Soft Computing for Problem Solving, December 28-30, 2012[C]. New Delhi: Springer, 2014.

\vspace{0.4cm}
\noindent {\color{red} \textbf{3．报告}}\\
{[1]} 中华人民共和国国务院新闻办公室．国防白皮书：中国武装力量的多样化运用[R/OL]．(2013-04-16)[2014-06-11]．\url{http://www.mod.gov.cn/affair/2013-04/16/content_4442839.htm}．\\
{[2]} World Health Organization．Factors regulating the immune response：report of WHO Scientific Group[R]．Geneva：WHO，1970．

\vspace{0.4cm}
\noindent {\color{red} \textbf{4．学位论文}}\\
{[1]} 吴云芳．面向中文信息处理的现代汉语并列结构研究[D/OL]．北京：北京大学，2003[2013-10-14]．\url{http://thesis.lib.pku.edu.cn/dlib/List.asp?lang=gb&type=Reader&DocGroupID=4&DocID=6328}．\\
{[2]} CALMS R B．Infrared spectroscopic studies on solid oxygen[D]．Berkeley：Univ．of California，1965．

\vspace{0.4cm}
\noindent {\color{red} \textbf{5．专利文献}}\\
{[1]} 张凯军．轨道火车及高速轨道火车紧急安全制动辅助装置：201220158825．2[P]．2012-04-05．\\
{[2]} 河北绿洲生态环境科技有限公司．一种荒漠化地区生态植被综合培植方法：01129210．5[P/OL]．2001-10-24[2002-05-28]．\url{http://211.152.9.47/sipoasp/zlijs/hyjs-yx-new.asp?recid=01129210.5&leixin=0}．\\
{[3]} KOSEKI A, MOMOSE H, KAWAHITO M, et al. Compiler: US828402[P/OL]. 2002-05-25[2002-05-28]. \url{http://FF&p=1&u=netahtml/PTO/search-bool.html&r=5&f=G&l=50&col=AND&d=PG01&sl=IBM.AS.&OS=AN/IBM&RS=AN/IBM}.

\vspace{0.4cm}
\noindent {\color{red} \textbf{6．标准文献}}\\
{[1]} 全国信息与文献标准化技术委员会．文献著录：第4部分 非书资料：GB/T 3792.4-2009[S]．北京：中国标准出版社，2010:3．

\vspace{0.4cm}
\noindent {\color{red} \textbf{7．专著中析出的文献}}\\
{[1]} 国家标准局信息分类编码研究所．世界各国和地区名称代码：GB/T 2659-1986[S]//全国文献工作标准化委员会．文献工作国家标准汇编：3．北京：中国标准出版社，1988：59-92．\\
{[2]} BUSECK P R, NORD G L, Jr, VEBLEN D R. Subsolidus phenomena in pyroxenes[M]//Pyroxenes. Washington, D.C.: Mineralogical Society of America, c1980: 117-211. \\
{[3]} FOURNEY M E. Advances in holographic photoelasticity[C]//Symposium on Applications of Holography in Mechanics, August 23-25, 1971, University of Southern California, Los Angeles, California. New York: ASME, c1971: 17-38.

\vspace{0.4cm}
\noindent {\color{red} \textbf{8．期刊中析出的文献}}\\
{[1]} 李炳穆．韩国图书馆法[J]．图书情报工作，2008，52（6）：6-21．\\
{[2]} 于潇，刘义，柴跃廷，等．互联网药品可信交易环境中主体资质审核备案模式[J]．清华大学学报(自然科学版)，2012，52(11)：1518-1523．\\
{[3]} DES MARAIS D J, STRAUSS H, SUMMONS R E, et al. Carbon isotope evidence for the stepwise oxidation of the Proterozoic environment[J]. Nature, 1992, 359: 605-609.

\vspace{0.4cm}
\noindent {\color{red} \textbf{9．报纸中析出的文献}}\\
{[1]} 丁文祥．数字革命与竞争国际化[N]．中国青年报，2000-11-20（15）．\\
{[2]} 张田勤．罪犯DNA库与生命伦理学计划[N]．大众科技报，2000-11-12（7）．

\vspace{0.4cm}
\noindent {\color{red} \textbf{10．电子资源（不包括电子专著、电子连续出版物、电子学位论文、电子专利）}}\\
{[1]} 萧钰．出版业信息化迈入快车道[EB/OL]．(2001-12-19)[2002-04-15]．\url{http://www.creader.com/news/20011219/200112190019.html}．\\
{[2]} 李强．化解医患矛盾需釜底抽薪[EB/OL]．(2012-05-03)[2013-03-25]．\url{http://wenku.baidu.com/view/47e4f206b52acfc789ebc92f.html}．\\
{[3]} Dublin core metadata element set: version 1.1[EB/OL]. (2012-06-14)[2014-06-11]. \url{http://dublincore.org/documents/dces/}．
}

\vspace{0.6cm}
\begin{table}[htbp]
  \centering
  \caption{文献类型、标识代码和电子资源载体、标识代码}
  \vspace{6pt}
  \fontsize{10.5pt}{14pt}\selectfont
  \begin{tabular}{|m{3.5cm}|m{2.5cm}|m{4.5cm}|m{2.5cm}|}
    \hline
    \multicolumn{2}{|c|}{\textbf{文献类型和标识代码}} & \multicolumn{2}{c|}{\textbf{电子资源载体和标识代码}} \\ \hline
    \centering 参考文献类型 & \centering 标识代码 & \centering 电子资源的载体类型 & \centering 标识代码 \tabularnewline \hline
    \centering 普通图书 & \centering M & \centering 磁带（magnetic tape） & \centering MT \tabularnewline \hline
    \centering 会议录 & \centering C & \centering 磁盘（disk） & \centering DK \tabularnewline \hline
    \centering 汇编 & \centering G & \centering 光盘（CD-ROM） & \centering CD \tabularnewline \hline
    \centering 报纸 & \centering N & \centering 联机网络（online） & \centering OL \tabularnewline \hline
    \centering 期刊 & \centering J & & \tabularnewline \hline
    \centering 学位论文 & \centering D & & \tabularnewline \hline
    \centering 报告 & \centering R & & \tabularnewline \hline
    \centering 标准 & \centering S & & \tabularnewline \hline
    \centering 专利 & \centering P & & \tabularnewline \hline
    \centering 数据库 & \centering DB & & \tabularnewline \hline
    \centering 计算机程序 & \centering CP & & \tabularnewline \hline
    \centering 电子公告 & \centering EB & & \tabularnewline \hline
    \centering 档案 & \centering A & & \tabularnewline \hline
    \centering 舆图 & \centering CM & & \tabularnewline \hline
    \centering 数据集 & \centering DS & & \tabularnewline \hline
    \centering 其他 & \centering Z & & \tabularnewline \hline
  \end{tabular}
\end{table}

% ==================== 8. 附录 ====================
\clearpage
\addcontentsline{toc}{chapter}{\texorpdfstring{附录 A $\times\times\times\times\times$}{附录 A XXXXX}}
\begin{center}
  {\heiti\bfseries\fontsize{16pt}{16pt}\selectfont 附录 A $\times\times\times\times\times$}
  \vspace{18pt}
\end{center}

……

% ==================== 9. 致谢 ====================
\clearpage
\addcontentsline{toc}{chapter}{\texorpdfstring{致\quad 谢}{致谢}}
\begin{center}
  {\heiti\bfseries\fontsize{16pt}{16pt}\selectfont 致\quad 谢}
  \vspace{18pt}
\end{center}

% 严格执行批注：仿宋 12 磅，行距 1.5 倍
\begin{spacing}{1.5}
{\fangsong\fontsize{12pt}{12pt}\selectfont
致谢应放置在参考文献、附录页后，另起页。致谢是作者对该论文的形成作出过贡献的组织或个人予以感谢的文字记载，语言要诚恳、恰当、简短。致谢对象应包括国家科学基金；资助研究工作的奖学金基金、合同单位、资助或支持的企业、组织或个人；协助完成研究工作和提供便利条件的组织或个人；在研究工作中提出建议和提供帮助的人。

字数一般不超过 1000 字。

\vspace{3.5cm}
\begin{flushright}
  学生姓名\qquad\qquad\quad \\[0.4cm]
  \dateunderline 于广州体育学院
\end{flushright}
}
\end{spacing}

% ==================== 10. 个人简历 ====================
\clearpage
\addcontentsline{toc}{chapter}{个人简历}
\begin{center}
  {\heiti\bfseries\fontsize{16pt}{16pt}\selectfont 个人简历}
  \vspace{18pt}
\end{center}

% 严格执行批注：仿宋 12 磅（英文 Times New Roman 12 磅），行距 1.5 倍，特定词组加粗
\begin{spacing}{1.5}
{\fangsong\fontsize{12pt}{12pt}\selectfont
\noindent \textbf{个人经历：}\\
$\times\times\times\times$，男，1975年8月生。\\
1997年7月毕业于$\times\times$大学体育教育专业，获学士学位。\\
1997年9月进入广州体育学院攻读硕士学位。

\vspace{0.4cm}
\noindent \textbf{在读期间发表论文：}\\
{[1]} $\times\times\times$.$\times\times\times\times\times$.$\times\times\times\times\times$, 2001, Vol.21 (3):70-74

\vspace{0.4cm}
……

\vspace{0.4cm}
\noindent \textbf{研究成果：}\\
{[1]} $\times\times$, $\times\times$, $\times\times$. $\times\times\times\times\times\times\times$研究. 2025

\vspace{0.4cm}
……
}
\end{spacing}

\end{document}